\documentclass[25pt, a0paper, landscape]{tikzposter}

% ---------- Packages ----------
\usepackage[T1]{fontenc}
\usepackage[utf8]{inputenc}
\usepackage{lmodern}
\usepackage{microtype}
\usepackage{amsmath,amssymb}
\usepackage{graphicx}
\usepackage{booktabs,array}
\usepackage{enumitem}
\usepackage{qrcode}
\usepackage{tikz}
\usetikzlibrary{calc,positioning}

% ---------- MindBridge Forest/Teal Green Palette ----------
\definecolor{MBGreen}{HTML}{0B5D4B}        % Deep forest/teal green (matches UI buttons/badges)
\definecolor{MBGreenDark}{HTML}{074235}    % Darker green accent
\definecolor{MBGreenLight}{HTML}{E8F3F0}   % Soft mint tint for callout boxes
\definecolor{MBGreenLine}{HTML}{9CC5BA}    % Muted sage border
\definecolor{MBText}{HTML}{182421}         % Dark charcoal-green text for high contrast
\definecolor{MBMuted}{HTML}{526660}
\definecolor{MBWhite}{HTML}{FFFFFF}

% ---------- Theme ----------
\usetheme{Simple}
\usecolorstyle{Default}
\usebackgroundstyle{Default}

\colorlet{backgroundcolor}{MBWhite}
\colorlet{framecolor}{MBGreenLine}
\colorlet{titlefgcolor}{white}
\colorlet{titlebgcolor}{MBGreen}
\colorlet{blocktitlebgcolor}{MBGreen}
\colorlet{blocktitlefgcolor}{white}
\colorlet{blockbodybgcolor}{white}
\colorlet{blockbodyfgcolor}{MBText}

\setlength{\parindent}{0pt}
\setlength{\parskip}{0.08em}
\linespread{0.96}
\setlist[itemize]{leftmargin=1.05em,itemsep=0.01em,topsep=0.01em,parsep=0pt,partopsep=0pt}

% ---------- Block appearance ----------
\defineblockstyle{MBBlock}{
    titlewidthscale=1, bodywidthscale=1, roundedcorners=16,
    titleleft, titleoffsetx=0mm, titleoffsety=0mm,
    bodyoffsetx=0mm, bodyoffsety=2.5mm,
    bodyverticalshift=2.5mm, linewidth=1.0mm
}{
    \draw[rounded corners=16pt, line width=1.1mm, color=MBGreenLine, fill=white]
      (blockbody.south west) rectangle (blockbody.north east);
    \ifBlockHasTitle
      \draw[rounded corners=16pt, line width=0mm, fill=MBGreen]
      (blocktitle.south west) rectangle (blocktitle.north east);
    \fi
}
\useblockstyle{MBBlock}

% ---------- Helpers ----------
\newcommand{\takeawaybox}[1]{%
  \begin{center}
  \fcolorbox{MBGreenLine}{MBGreenLight}{
    \parbox{0.94\linewidth}{\centering\bfseries\Large #1}
  }
  \end{center}
}

\newcommand{\posterimageplaceholder}[2][0.95\linewidth]{%
  \begin{center}
  \fcolorbox{MBGreenLine}{MBGreenLight}{
    \parbox[c][6.6cm][c]{#1}{\centering\bfseries\Large #2\\[0.35em]\normalsize Replace with a polished figure, screenshot, architecture diagram, chart, or hardware photo}
  }
  \end{center}
}

\newcommand{\metricbox}[2]{%
  \fcolorbox{MBGreenLine}{MBGreenLight}{
    \parbox[c][2.9cm][c]{0.29\linewidth}{\centering\bfseries\LARGE #1\\[0.18em]\normalsize #2}
  }
}

\newcommand{\projectmeta}[4]{%
  \begin{tabular}{>{\bfseries}p{0.24\linewidth}p{0.69\linewidth}}
    Team & #1\\[0.12em]
    Programme & #2\\[0.12em]
    Supervisors & #3\\[0.12em]
    Demo / QR & #4
  \end{tabular}
}

% ---------- Title (Wrapped cleanly so it never clips on the right) ----------
\title{\parbox{0.72\linewidth}{\centering MindBridge: Academic Workload Deconstruction\\and Burnout Mitigation Platform}}
\author{Tendai Lynn Mtakiwa \textbar\ Maunice Akaliza \textbar\ Yves Isite \textbar\ Orla Ishimwe \textbar\ Esther Konde \textbar\ Joel Hakizimana \textbar\ Muuse Muuse}
\institute{African Leadership University \textbar\ BSc. in Software Engineering Foundations Project \textbar\ Group Name: MindBridge}
\titlegraphic{\includegraphics[height=3.5cm]{alu-logo-white-full.png}}

\begin{document}
\maketitle

\vspace{0.05cm}

\begin{columns}

% ================= LEFT =================
\column{0.30}

% --- SECTION 1: MAUNICE AKALIZA ---
\block{1. Problem, Context, and Users}{
\begin{itemize}
  \item \textbf{Context \& Stress:} ALU students face heavy, overlapping deadlines; an initial survey found 83\% experienced moderate to very high academic stress.
  \item \textbf{Current Gap:} Canvas lists due dates but does not break assignments into manageable tasks.
  \item \textbf{Solution Direction:} MindBridge helps students organize assignments and plan their workload.
\end{itemize}
}

% --- SECTION 2: MAUNICE AKALIZA ---
\block{2. Objectives and Scope}{
\begin{itemize}
  \item \textbf{Objective 1:} Build MindBridge, a web platform that simplifies workload management and reduces last-minute pressure.
  \item \textbf{Scope:} Imports Canvas assignments, manages deadlines, and breaks rubrics into dated milestones for ALU students.
\end{itemize}
}

% --- SECTION 3: TENDAI MTAKIWA (COMPLETED) ---
\block{3. Project Metadata}{
\projectmeta
  {\textbf{Group MindBridge (Team 2):} Yves Isite (Tech Lead), Tendai Lynn Mtakiwa (Scrum Master \& QA), Muuse Muuse (Product Owner \& QA), Hakizimana Baraka Joel (Full Stack \& QA), Orla Ishimwe (Full Stack), Esther Konde (System Architect), Akaliza Alizee Maunice (PM \& QA)}
  {BSc. in Software Engineering Foundations Project (ALU Rwanda)}
  {Foundations Project Facilitators \& Faculty Mentors}
  {\texttt{github.com/IsiteYves/mindbridge}}

\vspace{0.25em}
\begin{center}
\qrcode[height=3.4cm]{https://github.com/IsiteYves/mindbridge}
\end{center}
}

% --- SECTION 4: JOEL & ESTHER ---
\block{4. Constraints and Design Decisions}{
Posters are stronger when they show real engineering trade-offs.

\begin{itemize}
  \item Connectivity, power, budget, or hardware constraints
  \item Data quality, privacy, security, or compliance issues
  \item Why you chose this architecture, model, stack, or device setup
  \item Main trade-off: what you optimized for, and what you sacrificed
\end{itemize}
}

% ================= MIDDLE =================
\column{0.40}

% --- SECTION 5: ORLA, JOEL & ESTHER ---
\block{5. Proposed Solution and System Overview}{
\takeawaybox{End-to-End Flow: Connect Canvas $\rightarrow$ Break Rubrics into Milestones $\rightarrow$ Balance Daily Workload}

\begin{minipage}[t]{0.48\linewidth}
\textbf{Solution summary}

MindBridge helps ALU students turn a heavy assignment load into a plan they can finish. A student connects Canvas once, and MindBridge imports the assignments and breaks each rubric into dated milestones that the student can edit and reorder. A workload heatmap shows where work piles up, a focus timer and daily goal keep each day small, and reminder emails arrive before deadlines. MindBridge supports planning; it does not diagnose or treat mental health.
\end{minipage}
\hfill
\begin{minipage}[t]{0.48\linewidth}
\textbf{System architecture}
\begin{itemize}
  \item \textbf{Frontend:} React 19, TypeScript, Vite, \& React Router UI
  \item \textbf{Backend:} Flask 3.1 REST API, SQLAlchemy ORM, PyJWT auth, \& Fernet token encryption
  \item \textbf{Database \& AI:} PostgreSQL database \& optional LangGraph plan-review agent
  \item \textbf{Integrations:} Canvas LMS REST API (with manual import fallback) \& email reminders
\end{itemize}
\end{minipage}

\vspace{0.4em}
\begin{center}
\begin{minipage}[t]{0.23\linewidth}\centering
  \includegraphics[width=\linewidth]{figures/ui-canvas.png}\\[0.2em]
  {\small\textbf{1. Connect Canvas once}}
\end{minipage}\hfill
\begin{minipage}[t]{0.23\linewidth}\centering
  \includegraphics[width=\linewidth]{figures/ui-assignments.png}\\[0.2em]
  {\small\textbf{2. Import assignments}}
\end{minipage}\hfill
\begin{minipage}[t]{0.23\linewidth}\centering
  \includegraphics[width=\linewidth]{figures/ui-milestones.png}\\[0.2em]
  {\small\textbf{3. Rubric to milestones}}
\end{minipage}\hfill
\begin{minipage}[t]{0.23\linewidth}\centering
  \includegraphics[width=\linewidth]{figures/ui-heatmap.png}\\[0.2em]
  {\small\textbf{4. Track workload heatmap}}
\end{minipage}
\end{center}
}

% --- SECTION 6: YVES ISITE ---
\block{6. Methodology and Build Process}{
\begin{minipage}[t]{0.48\linewidth}
\textbf{Research / discovery}
\begin{itemize}
  \item User interviews, observation, or precedent review
  \item Requirement gathering and feasibility checks
  \item Risk, ethics, or context analysis
\end{itemize}
\end{minipage}
\hfill
\begin{minipage}[t]{0.48\linewidth}
\textbf{Engineering / implementation}
\begin{itemize}
  \item Architecture and design choices
  \item Prototype iterations and testing
  \item Deployment, simulation, or demo setup
\end{itemize}
\end{minipage}

\vspace{0.4em}
\begin{center}
\begin{tabular}{p{0.16\linewidth}p{0.76\linewidth}}
\toprule
\textbf{Stage} & \textbf{What to show briefly} \\
\midrule
Discover & users, needs, constraints, baseline findings \\
Design & flows, architecture, rationale, sketches \\
Build & key components, integrations, implementation choices \\
Test & usability, speed, accuracy, reliability, feedback \\
Iterate & what changed after evidence or testing \\
\bottomrule
\end{tabular}
\end{center}
}

% --- SECTION 7: JOEL HAKIZIMANA ---
\block{7. Results, Evaluation, and Demo Evidence}{
\begin{center}
\metricbox{95\%}{Task success}
\hfill
\metricbox{1.8s}{Avg. response}
\hfill
\metricbox{12}{Pilot users}
\end{center}

\vspace{0.5em}
\begin{minipage}[t]{0.48\linewidth}
\begin{itemize}
  \item Show 2--4 quantitative results only
  \item Include one user or field insight if relevant
  \item Compare against a baseline where possible
\end{itemize}
\begin{center}\fcolorbox{MBGreenLine}{MBGreenLight}{\parbox[c][3.7cm][c]{0.95\linewidth}{\centering\bfseries\Large Chart, benchmark, confusion matrix, or evaluation table}}\end{center}
\end{minipage}
\hfill
\begin{minipage}[t]{0.48\linewidth}
\begin{itemize}
  \item What can be demonstrated live today?
  \item Which workflow is fully working?
  \item What remains simulated or partial?
\end{itemize}
\begin{center}\fcolorbox{MBGreenLine}{MBGreenLight}{\parbox[c][3.7cm][c]{0.95\linewidth}{\centering\bfseries\Large Demo storyboard, product screenshot, or device photo}}\end{center}
\end{minipage}
}

% ================= RIGHT =================
\column{0.30}

% --- SECTION 8: MUUSE MUUSE ---
\block{8. Impact, Contribution, and Learning}{
Frame the contribution at three levels.

\begin{itemize}
  \item \textbf{Technical:} What did you build, optimize, or integrate?
  \item \textbf{User:} What becomes easier, faster, safer, cheaper, or more inclusive?
  \item \textbf{Research:} What did you learn, validate, or discover?
  \item \textbf{Mission:} Why does this matter in an ALU / African context?
\end{itemize}

\takeawaybox{Do not just show that it works. Show why it matters.}
}

% --- SECTION 9: MUUSE MUUSE ---
\block{9. Limitations and Next Steps}{
Strong projects acknowledge what is still unfinished.

\begin{itemize}
  \item Limitations in data, hardware, testing scope, or deployment
  \item Assumptions and known failure points
  \item Immediate next iteration toward real-world use
  \item Long-term future work or scaling path
\end{itemize}
}

% --- SECTION 10: MUUSE MUUSE ---
\block{10. References and Acknowledgements}{
\textbf{References}

\begin{itemize}
  \item Author, A. (Year). Title. Journal or Conference.
  \item Author, B. (Year). Title. Source.
  \item Author, C. (Year). Title. Source.
\end{itemize}

\textbf{Acknowledgements}

Thank supervisors, labs, partners, participants, funders, or collaborators.
}

\end{columns}

\end{document}